\section{Objetivos}

\subsection{Objetivo general}

Implementar el backend completo del sistema, desarrollando los servicios, repositorios, 
DTOs y controladores de los seis módulos funcionales (Usuarios, Autenticación, Catálogo, 
Productos, Dispositivos y Préstamos), exponiendo la API REST documentada 
que sirva como base para el desarrollo del frontend en los hitos siguientes.

\subsection{Objetivos específicos}

\begin{itemize}
	\item Documentar el uso de Alembic como herramienta de migración iterativa durante 
	el desarrollo, describiendo los ajustes realizados al esquema de base de datos 
	conforme evolucionaron los requerimientos del sistema.
	
	\item Implementar el módulo de Usuarios con sus métodos personalizados de creación 
	y actualización vinculados a roles, y con guards de acceso que diferencien usuarios 
	activos de administradores.
	
	\item Implementar el módulo de Autenticación, centralizando la emisión y validación de tokens 
	JWT para todas las rutas protegidas del sistema.
	
	\item Implementar el módulo de Catálogo para las 
	entidades de referencia del sistema: \texttt{roles}, \texttt{brands}, 
	\texttt{models}, \texttt{categories}, \texttt{statuses} y \texttt{ubications}.
	
	\item Implementar el módulo de Productos con sus campos de relación opcionales como \texttt{brand\_id}, 
	\texttt{model\_id} y \texttt{category\_id}.
	
	\item Implementar el módulo de Dispositivos con métodos personalizados de consulta, 
	filtrado y cambio de estado, incorporando un registro automático de historial.
	
	\item Implementar el módulo de Préstamos con el flujo completo de ciclo de vida 
	de una solicitud: creación, aprobación o rechazo, entrega y devolución, 
	actualizando los estados de los dispositivos involucrados.
\end{itemize}